\section{Availability and Reproducibility}

Epistract is released as open-source software under the MIT license and is available on GitHub. The system is installable as a Claude Code plugin in a single command, requiring no manual dependency resolution beyond the Claude Code runtime itself. On first invocation, the setup command automatically installs the Python extraction pipeline, the sift-kg graph construction library, and all required dependencies into an isolated virtual environment.

The runtime requirements are minimal: Claude Code as the execution environment, Python 3.11 or later, and the sift-kg library for graph construction. Two optional dependencies extend the pipeline's validation capabilities: RDKit enables SMILES string validation and molecular property computation for chemical entities, and Biopython enables sequence format validation for protein and nucleic acid entities. Both are installed automatically when available but degrade gracefully when absent, ensuring the core pipeline functions on any system that meets the base requirements.

All test artifacts are committed to the repository and available for inspection. The five validation corpora comprise 77 PubMed abstracts in plain text format, organized by therapeutic domain. For each scenario, the repository contains the raw extraction outputs in JSON format, the constructed graph data, and evidence screenshots of the final knowledge graph visualizations. This means any reviewer can re-run any scenario against the committed corpus and compare the resulting graph against the committed baseline, providing a concrete and verifiable reproducibility guarantee rather than a procedural one.

The plugin architecture provides an additional reproducibility benefit: because epistract is distributed as a Claude Code plugin rather than a standalone application, updates to the extraction pipeline, domain schema, and validation logic are propagated automatically when the plugin is updated. Users always run the current version of the system without manual upgrade procedures, and the committed test baselines serve as regression anchors across versions.

\subsection{Open-Source Vision and Community Contribution}

Epistract is open-sourced not merely for transparency but with the explicit goal of community adoption and collaborative enrichment. The domain schema (\texttt{domain.yaml}) is designed to be extensible: contributors can add new entity types, relation types, and ontology mappings to cover therapeutic areas beyond the five validated here. The extraction skill prompt can be adapted for adjacent scientific domains --- veterinary medicine, agricultural biotechnology, materials science --- where the same pattern of structured knowledge extraction from literature applies but the entity vocabulary differs.

I encourage contributions in several forms: new evaluation scenarios with curated PubMed corpora that test the schema against additional drug discovery domains; schema extensions that introduce entity types or relation types for specialized subdisciplines (e.g., pharmacokinetics, toxicogenomics); additional export formats for integration with platforms such as Neo4j, Cytoscape, or institutional knowledge management systems; and external database connectors that enrich extracted entities with data from PubChem, ChEMBL, UniProt, or ClinicalTrials.gov. The broader aspiration is that epistract becomes a shared instrument --- a community-maintained bridge between scientific literature and structured knowledge --- advancing the adoption of AI-assisted workflows in pharmaceutical and biomedical research.

\subsection{Limitations and Responsible Use}

Epistract is a research tool, not a clinical or regulatory decision-making system. The knowledge graphs it produces are derived from LLM extraction, which is inherently probabilistic: entities may be misclassified, relations may be inferred incorrectly, and confidence scores reflect model certainty rather than ground truth. All outputs should be treated as hypotheses that require human expert review and validation before informing any scientific, clinical, or regulatory decision. The system is not a substitute for systematic literature review, pharmacovigilance processes, or evidence-based clinical practice. Users in pharmaceutical and healthcare settings should apply appropriate institutional review and domain expert oversight before acting on any information surfaced by the pipeline.
